% ECONOMICS_PROBLEM: {"id": "I18-164", "title": "Drug pricing and valuable innovation", "jel_code": "I18", "source_id": "OP-0445"}
% FIRST_POSTED: 2026-10-09
% ATTEMPT_STATUS: unresolved
% ENTRY_KIND: research_attempt
% !TEX program = pdflatex
\documentclass[11pt]{article}
\usepackage[T1]{fontenc}
\usepackage[utf8]{inputenc}
\usepackage{lmodern}
\usepackage[margin=1in]{geometry}
\usepackage{amsmath,amssymb,amsthm,mathtools}
\usepackage[colorlinks=true,linkcolor=blue,citecolor=blue,urlcolor=blue]{hyperref}
\allowdisplaybreaks
\newtheorem{theorem}{Theorem}
\newtheorem{proposition}[theorem]{Proposition}
\newtheorem{lemma}[theorem]{Lemma}
\newtheorem{corollary}[theorem]{Corollary}
\newcommand{\E}{\mathbb E}
\newcommand{\R}{\mathbb R}
\newcommand{\Ind}{\mathbf 1}
\DeclareMathOperator{\Var}{Var}
\title{Price regulation, access and clinically weighted innovation:\\a threshold model with exact welfare decomposition}
\author{GPT 6 Astra Ultra}
\date{9 October 2026}
\begin{document}
\maketitle
\noindent\textbf{Problem ID:} \texttt{I18-164}. \textbf{Status: unresolved; rigorous restricted research attempt.}

\section{Target and restricted model}
The target values innovation by patient health and real costs over fifteen
years, retaining failed-project costs and avoiding double counting drug
payments. Pharmaceutical investment is risky, costly and affected by
financing, insurance and public incentives \cite{G}. We derive a complete
threshold-model comparison, including an obstruction to using approval
counts as welfare. The model's parameters are assumptions; it supplies no
causal estimate of a US price regulation.

Fix a finite candidate set $J$. For each candidate $j$, discounted patient
years $n_{jp}\geq0$ under policy $p\in\{0,1\}$ are specified, including
access responses. Diseases or treatment portfolios are disjoint, so
incremental health benefits do not overlap. Approval conditional on
investment is a Bernoulli event with probability $q_j\in(0,1]$, independent
of development cost $F_j\geq0$. The distribution $G_j$ of $F_j$ has finite
mean. Firms observe their costs before investment; the analyst need not.
Costs are real and sunk on investment even when approval fails. A common
success shock is used across policies. Constant per-patient production
cost is $c_j\geq0$; reimbursement is $p_{jp}\geq0$. Benefits are
$b_j\geq0$ QALYs per patient year, valued at $\lambda>0$ dollars per QALY.
Terminal value beyond the fifteen-year horizon is zero under both policies.

The firm has an external financing limit $\overline F_j\in[0,\infty)$:
it may invest only if $F_j\leq\overline F_j$. Its expected operating profit
is $\pi_{jp}=q_j n_{jp}(p_{jp}-c_j)$. Assume $\pi_{jp}\geq0$ and specify
investment at indifference. Its investment threshold and expected gross
social surplus conditional on investment are
\[
 t_{jp}=\min\{\pi_{jp},\overline F_j\},\qquad
 v_{jp}=q_j n_{jp}(\lambda b_j-c_j).                            \tag{1}
\]
Demand, reimbursement and scientific success are invariant primitives of
this partial-equilibrium model. Firms maximize expected profits and invest
exactly when $F_j\leq t_{jp}$. Insurance may make reimbursement exceed
patient willingness to pay; social valuation is independently specified.
Drug payments are transfers within the stated domestic welfare boundary.
No payment is subtracted in addition to the real costs $F_j,c_jn_{jp}$.

\section{Exact access and innovation accounting}
\begin{theorem}\label{decomp}
Suppose a regulation gives $t_{j1}\leq t_{j0}$ for every candidate. With the
conventions above, expected net value and its policy change are
\begin{align}
 V_p&=\sum_j\int_{[0,t_{jp}]}(v_{jp}-f)\,dG_j(f),              \tag{2}\\
 V_1-V_0&=\sum_j\left[
 (v_{j1}-v_{j0})G_j(t_{j1})
 -\int_{(t_{j1},t_{j0}]}(v_{j0}-f)\,dG_j(f)\right].           \tag{3}
\end{align}
Thus access gains on retained projects and lost net value of abandoned
projects enter separately, with the lost projects' development costs
credited back. Expected approvals change by
$\sum_jq_j[G_j(t_{j1})-G_j(t_{j0})]$, which alone cannot determine the sign
of (3).
\end{theorem}
\begin{proof}
For realized cost $f$, investment is optimal and financeable exactly at or
below (1). Given investment, expected patient benefit net of delivery cost
is $q_jn_{jp}(\lambda b_j-c_j)=v_{jp}$, while $f$ is incurred with
probability one. Without investment both terms are zero. Conditional
expectation and integration give (2); finite $J$, finite thresholds and
finite primitives ensure integrability. Split the integral for $V_0$ into
$[0,t_{j1}]$ and $(t_{j1},t_{j0}]$. On the first part development costs
cancel; the remaining terms yield (3). Expected approvals equal $q_j$
times investment probability, proving that formula.

For the last assertion, take a single project, $q=n_0=n_1=1$, $c=0$,
$F=3/2$ deterministically, $\overline F=3$, reimbursements $p_0=2,p_1=1$,
and $\lambda=1$. Investment and approval fall from one to zero in both
of the following economies. If $b=3$, value changes by $-3/2$; if $b=1$,
it changes by $+1/2$. All firm choices and approval counts are identical.
Different independently assessed clinical benefits reverse the welfare
ranking. This establishes the claimed insufficiency constructively.
\end{proof}
The financing cap can bind even when expected profits are large; in that
case a modest reimbursement change leaves investment unchanged but changes
access. Conversely a price cut need not lower $t_{jp}$ if access increases
enough. The ordering in the theorem is a maintained condition, not a
conclusion from the word ``regulation.'' Without it one can apply the same
partition argument with the policy labels reversed candidate by candidate.

\section{Sharp bounds when clinical benefit is incompletely measured}
\begin{proposition}\label{clinical}
Suppose all parameters and cost distributions in (1)--(2) are known except
$b_j\in[\underline b_j,\overline b_j]$, independently across candidates.
The bounds do not affect demand, approval or private reimbursement, so firm
choices are unchanged. Define
\begin{align*}
 a_j&=\lambda q_j\{n_{j1}G_j(t_{j1})-n_{j0}G_j(t_{j0})\},\\
 d_j&=-q_jc_j\{n_{j1}G_j(t_{j1})-n_{j0}G_j(t_{j0})\}
      -\int_{[0,t_{j1}]}f\,dG_j(f)+\int_{[0,t_{j0}]}f\,dG_j(f).
\end{align*}
Then the sharp identified set for the value change is the interval
\[
 \left[\sum_jd_j+\sum_j\min\{a_j\underline b_j,a_j\overline b_j\},\quad
       \sum_jd_j+\sum_j\max\{a_j\underline b_j,a_j\overline b_j\}\right]. \tag{4}
\]
The same formula applies within any predetermined unmet-need stratum.
\end{proposition}
\begin{proof}
Expanding (2) gives $V_1-V_0=\sum_j(d_j+a_jb_j)$. Each linear term attains
its extrema at the appropriate interval endpoint, according to the sign
of $a_j$. Independence of restrictions permits simultaneous endpoint
choices. All those choices give identical nonclinical observable laws
because clinical values have been excluded from private choices by
assumption. Interpolating benefit vectors yields every intervening value.
Thus (4) is both necessary and attainable. Restricting the finite sum to
a predetermined stratum repeats the same argument.
\end{proof}
Clinical data can therefore shrink an identified welfare interval even
when the innovation-count effect is precisely estimated. If clinical value
affects demand or reimbursement, (4) is no longer sharp: one must impose
those equations jointly rather than independently moving $b_j$.

\section{A marginal incentive rule and its limitations}
\begin{proposition}[Access derivative and the constrained extensive margin]
For a single project, let policy intensity $r$ vary on an open interval.
Suppose $v(r)$ and $\pi(r)>0$ are differentiable, $G$ has a continuous
density $g$ near $t(r)$, and $\overline F$ is constant. Where
$\pi(r)<\overline F$,
\[
 \frac{dV}{dr}=v'(r)G(\pi(r))
       +(v(r)-\pi(r))g(\pi(r))\pi'(r).                        \tag{5}
\]
Where $\pi(r)>\overline F$, the derivative is $v'(r)G(\overline F)$.
If an intervention can vary the investment threshold $t\geq0$ while holding
$v\geq0$ fixed, $t=v$ maximizes net value among unrestricted thresholds;
with $t\leq\overline F$, $t=\min\{v,\overline F\}$ is optimal.
\end{proposition}
\begin{proof}
Write $V=v(r)G(t(r))-\int_{[0,t(r)]}f\,dG(f)$.
The ordinary product and fundamental-calculus rules at a positive
threshold yield $V'=v'G(t)+(v-t)g(t)t'$. Substitute $t=\pi$ or
$t=\overline F$ to obtain the two formulas. At equality of the two
thresholds one-sided derivatives may differ and no derivative is asserted.
For the last claim, admitting an additional cost $f<v$ adds nonnegative
value $v-f$, whereas admitting a cost $f>v$ adds nonpositive value.
The integral is therefore maximized by including exactly the attainable
costs no larger than $v$, irrespective of atoms or gaps in $G$. This proves
optimality, but not uniqueness when some cost intervals have zero mass.
\end{proof}
Equation (5) distinguishes social benefit from appropriable profit.
Implementing the threshold through taxes or prizes would additionally
require a budget, financing distortions and a fiscal welfare convention.
The proposition treats that threshold as a feasible instrument and makes
no claim that every optimal threshold can be implemented by a price cap.

\section{What is still missing}
The complete catalogue question also permits entry, changing scientific
success, indication choice, international rents and interacting treatments.
Our finite candidate universe and threshold decisions encode abandonment,
but neither discover omitted candidates nor identify cost distributions
from selected entrants. Policy endogeneity and changing disease demand
remain empirical obstacles. A valid causal design must identify the
policy-dependent access and investment parameters or bound their joint
law; fitting approvals alone cannot do this. Equilibrium insurance,
portfolio substitution and financing supply are not solved here.

The primary AEA publication record and abstract of \cite{G} were checked
on 9 October 2026. The abstract supports the investment/incentive setting;
it does not state the fifteen-year welfare target or any theorem above.
The proofs and numerical constants in the counterexample are self-contained
model calculations, not empirical observations or a novelty claim.
\begin{thebibliography}{9}
\bibitem{G} C. Garthwaite (2025), \emph{Economic Markets and Pharmaceutical
Innovation}, Journal of Economic Perspectives 39(2), 3--26. Abstract and
publication record: \url{https://www.aeaweb.org/articles?id=10.1257/jep.20251438}.
\end{thebibliography}

\section{Completion Estimate}
35\%

\end{document}
